\section*{Chapitre \ref{ch:g9:inside-the-atom} --- À l'intérieur de l'atome}

\begin{solution}{exo:g9:inside-the-atom:1}
Dans le \omterm{def:g9:inside-the-atom:nucleus}{noyau}~: des \omterm{def:g9:inside-the-atom:proton}{protons} (positifs) et des \omterm{def:g9:inside-the-atom:proton}{neutrons} (sans charge).
Autour~: des \omterm{def:g9:inside-the-atom:nucleus}{électrons} (négatifs).
\end{solution}

\begin{solution}{exo:g9:inside-the-atom:2}
\ce{^{16}_{8}O}~: 8 \omterm{def:g9:inside-the-atom:proton}{protons}, 8 \omterm{def:g9:inside-the-atom:proton}{neutrons}, 8 \omterm{def:g9:inside-the-atom:nucleus}{électrons}.
\ce{^{27}_{13}Al}~: 13, 14, 13. \ce{^{40}_{20}Ca}~: 20, 20, 20.
\end{solution}

\begin{solution}{exo:g9:inside-the-atom:3}
Il possède autant d'\omterm{def:g9:inside-the-atom:nucleus}{électrons} (charge $-e$ chacun) que de \omterm{def:g9:inside-the-atom:proton}{protons}
(charge $+e$ chacun)~: les charges se compensent.
\end{solution}

\begin{solution}{exo:g9:inside-the-atom:4}
$A = 26 + 30 = 56$~: \ce{^{56}_{26}Fe}.
\end{solution}

\begin{solution}{exo:g9:inside-the-atom:5}
Oui~: les deux ont $Z = 17$, ce sont donc deux \omterm{def:g7:atoms-and-molecules:atom}{atomes} de chlore. Ils ont
des \omterm{def:g9:inside-the-atom:atomic-number}{nombres de masse} différents (18 et 20 \omterm{def:g9:inside-the-atom:proton}{neutrons})~: ce sont des
\omterm{def:g9:inside-the-atom:isotope}{isotopes}.
\end{solution}

\begin{solution}{exo:g9:inside-the-atom:6}
Non~: ils ont des \omterm{def:g9:inside-the-atom:atomic-number}{numéros atomiques} différents (6 et 7), ce sont donc
des éléments différents, le carbone et l'azote. Des \omterm{def:g9:inside-the-atom:isotope}{isotopes} ont le même
$Z$.
\end{solution}

\begin{solution}{exo:g9:inside-the-atom:7}
$56 \times 1.67 \times 10^{-27} \approx \qty{9.35e-26}{kg}$.
\end{solution}

\begin{solution}{exo:g9:inside-the-atom:8}
Un \omterm{def:g7:atoms-and-molecules:atom}{atome} est surtout fait de vide~: le \omterm{def:g9:inside-the-atom:nucleus}{noyau} est minuscule, si bien que
la plupart des particules ne rencontrent aucun \omterm{def:g9:inside-the-atom:nucleus}{noyau} sur leur chemin et
traversent tout droit.
\end{solution}

\begin{solution}{exo:g9:inside-the-atom:9}
$10^{-10} / 10^{-15} = 10^{5}$~: cent mille fois plus petit.
\end{solution}

\begin{solution}{exo:g9:inside-the-atom:10}
La chimie dépend des \omterm{def:g9:inside-the-atom:nucleus}{électrons}, et les \omterm{def:g9:inside-the-atom:isotope}{isotopes} d'un élément ont le même
nombre d'\omterm{def:g9:inside-the-atom:nucleus}{électrons} (le même $Z$).
\end{solution}

\begin{solution}{exo:g9:inside-the-atom:11}
$26 \times 9.11 \times 10^{-31} = \qty{2.37e-29}{kg}$. Fraction~:
$2.37 \times 10^{-29} / 9.35 \times 10^{-26} \approx 2.5 \times
10^{-4}$, soit environ \qty{0.025}{\percent} de la masse.
\end{solution}

\begin{solution}{exo:g9:inside-the-atom:12}
6915 \omterm{def:g7:atoms-and-molecules:atom}{atomes} de cuivre~63 et 3085 \omterm{def:g7:atoms-and-molecules:atom}{atomes} de cuivre~65. Moyenne~:
$(6915 \times 63 + 3085 \times 65)/10\,000 = 63.6$ \omterm{def:g9:inside-the-atom:proton}{nucléons}.
\end{solution}

\begin{solution}{pb:g9:inside-the-atom:1}
\textbf{1.} \ce{^{35}_{17}Cl} et \ce{^{37}_{17}Cl}.

\textbf{2.} Chlore~35~: 17 \omterm{def:g9:inside-the-atom:proton}{protons}, 18 \omterm{def:g9:inside-the-atom:proton}{neutrons}, 17 \omterm{def:g9:inside-the-atom:nucleus}{électrons}.
Chlore~37~: 17 \omterm{def:g9:inside-the-atom:proton}{protons}, 20 \omterm{def:g9:inside-the-atom:proton}{neutrons}, 17 \omterm{def:g9:inside-the-atom:nucleus}{électrons}.

\textbf{3.} Les deux ont 17 \omterm{def:g9:inside-the-atom:proton}{protons}~: ils appartiennent tous deux à
l'élément chlore~; ils ne diffèrent que par leur nombre de \omterm{def:g9:inside-the-atom:proton}{neutrons}~: ce
sont des \omterm{def:g9:inside-the-atom:isotope}{isotopes}.

\textbf{4.} Oui~: la chimie dépend des \omterm{def:g9:inside-the-atom:nucleus}{électrons}, et les deux en ont 17.

\textbf{5.} $35 \times 1.67 \times 10^{-27} = \qty{5.845e-26}{kg}$.

\textbf{6.} $17 \times 9.11 \times 10^{-31} = \qty{1.55e-29}{kg}$~:
environ 4000 fois moins que le \omterm{def:g9:inside-the-atom:nucleus}{noyau}, négligeable.

\textbf{7.} $37 \times 1.67 \times 10^{-27} = \qty{6.179e-26}{kg}$.

\textbf{8.} $\qty{1}{g} = \qty{e-3}{kg}$~; $10^{-3} / 5.845 \times
10^{-26} \approx 1.71 \times 10^{22}$ \omterm{def:g7:atoms-and-molecules:atom}{atomes}.

\textbf{9.} \qty{75.8}{\percent} et \qty{24.2}{\percent}.

\textbf{10.} $758 \times 5.845 \times 10^{-26} = \qty{4.430e-23}{kg}$~;
$242 \times 6.179 \times 10^{-26} = \qty{1.495e-23}{kg}$.

\textbf{11.} $4.430 \times 10^{-23} + 1.495 \times 10^{-23} =
\qty{5.925e-23}{kg}$.

\textbf{12.} $5.925 \times 10^{-23} / 1000 \approx
\qty{5.93e-26}{kg}$~: c'est la masse moyenne d'un \omterm{def:g7:atoms-and-molecules:atom}{atome} de chlore.
\end{solution}
